\centerline{\bf \Huge Всероссная тренировка}

\problem{1}
Дано нечётное натуральное число $n>1$. На доске записаны числа

$$
n, n+1, n+2, \ldots, 2 n-1 .
$$


Докажите, что можно стереть одно из них так, чтобы сумма оставшихся чисел не делилась ни на одно из оставшихся чисел.

\problem{2}
Дано натуральное $n>1$. Найдите наибольшее вещественное $\alpha$, такое что для всех положительных чисел $a_1, \ldots, a_n, b_1, \ldots, b_n$, удовлетворяющих условию $a_1+\ldots+a_n=b_1+\ldots+b_n=1$, выполнено неравенство

$$
\frac{a_1^2}{a_1^3+b_1^5}+\ldots+\frac{a_n^2}{a_n^3+b_n^5}>\alpha .
$$


\problem{3}
Дано натуральное число $n$. Двое играют в такую игру. Изначально на столе лежит $2 n+1$ рублёвых монет, а на доске написано число $a=1$. Игроки ходят по очереди. Каждым ходом игрок либо не меняет, либо увеличивает на 2 число $a$, после чего забирает $a$ монет со стола. Игра заканчивается, когда на столе остаётся меньше, чем $a$ монет; выигрывает тот, у кого в этот момент монет больше, а если монет у игроков поровну, игра заканчивается вничью. Каким будет результат игры при правильной игре обоих игроков?

\problem{4}
В треугольнике $A B C$ отметили центр вписанной окружности $I$ и провели описанную окружность $\omega$. Вписанная окружность треугольника $A B C$ касается $B C$ в точке $D$. Перпендикуляр, восставленный в точке $I$ к прямой $A I$, перескает прямые $A B$ и $A C$ в точках $E$ и $F$ соответственно. Описанная окружность треугольника $A E F$ пересекает $\omega$ и прямую $A I$ в точках $G$ и $H$. Касательная в точке $G$ к $\omega$ пересекает $B C$ в точке $J$. Прямая $A J$ пересекается с $\omega$ в точке $K$. Докажите, что описанные окружности треугольников $D J K$ и $G I H$ касаются.